\documentclass[11pt]{article}
\usepackage[a4paper,margin=28mm]{geometry}
\usepackage{amsmath,amssymb,amsthm}
\usepackage[colorlinks=true,linkcolor=blue,urlcolor=blue,citecolor=blue]{hyperref}
\newtheorem*{theorem*}{Theorem}
\newtheorem*{corollary*}{Corollary}
\newenvironment{statement}[1]{\par\medskip\noindent\phantomsection\textbf{#1}\ }{\par\medskip}
\setlength{\emergencystretch}{3em}
\newcommand{\Q}{\mathbb Q}
\newcommand{\Z}{\mathbb Z}
\newcommand{\F}{\mathbb F}
\DeclareMathOperator{\End}{End}
\DeclareMathOperator{\Gal}{Gal}
\DeclareMathOperator{\Res}{Res}
\DeclareMathOperator{\disc}{disc}
\title{Rational \(5\)-torsion and the \(2\)-primary torsion of geometrically split abelian surfaces}
\author{}
\date{14 November 2026}

\begin{document}
\maketitle

\begin{theorem*}
Let \(A/\Q\) be an abelian surface that is simple over \(\Q\) and geometrically split. If \(A(\Q)\)
has a point of order \(5\), then
\[\#A(\Q)[2^\infty]\le4 .\]
\end{theorem*}

\begin{corollary*}
No genus-\(2\) Jacobian over \(\Q\) that is geometrically split and simple over \(\Q\) has a rational
point of order \(5\) and at least \(8\) rational points of \(2\)-power order. In particular
\href{https://genus-2-torsion.github.io/group.html?g=40}{\([40]\)} and
\href{https://genus-2-torsion.github.io/group.html?g=2.20}{\([2,20]\)} do not occur as the rational
torsion subgroup of such a Jacobian. Neither do \([2,2,10]\), \([2,2,20]\) and \([2,2,2,10]\), which
the census proof \cite{P2210} had already excluded.
\end{corollary*}

Here ``geometrically split'' means that \(A_{\bar\Q}\) is isogenous to a product of two elliptic
curves. We impose no restriction on the degree of this isogeny, on the field of definition of the
geometric endomorphisms or on the integral endomorphism ring, and we assume no polarization. The
bound \(4\) is attained, by Jacobians with rational torsion subgroup \([2,10]\)
(Remark~4). The theorem strengthens the result of \cite{P2210} (a point of order \(5\)
forces rational \(2\)-rank at most \(2\), for principally polarized surfaces).

\emph{Outline.} We follow the case distinction of \cite{P2210}. In the CM case, and when the field of
definition of the geometric endomorphisms is not quadratic, the arguments of \cite{P2210} need only
\(40\mid\#\widetilde A(\F_3)\) or no information about the \(2\)-primary part at all
(Section~\ref{other}). In the remaining cases \(A\) contains, over a quadratic field \(K\), an elliptic
curve \(E\) with \(E^\sigma\ne E\). In \cite{P2210} a principal polarization and rational \(2\)-rank
\(3\) were used to find an elliptic curve in \(A_K\) with full \(K\)-rational \(2\)-torsion. Here a
curve with full \(K\)-rational \(2\)-torsion and a \(K\)-point of order \(5\), \(K\)-isogenous to \(E\),
comes from \(\#A(\Q)[2^\infty]\ge8\) alone. The exact sequence
used for \([2,2,14]\) \cite[Proposition~2]{P2214} puts \(2\)-primary torsion of order at least \(4\) on
\(E\) or on \(A_K/E^\sigma\). By Kenku--Momose no elliptic curve over a quadratic field has a point of
order \(20\), so with the point of order \(5\) this torsion is the full \(2\)-torsion of a curve in the
isogeny class (Lemma~9). From there the analysis of quadratic points on \(X_1(2,10)\) in
\cite{P2210} goes through with \(8\) in place of \(40\) (Section~\ref{quadratic}).

\section{Conventions and inputs}\label{inputs}

We use the conventions of \cite[Section~1]{P2210}. For a number field \(k\), \(G_k=\Gal(\bar\Q/k)\).
For a quadratic field \(K\), \(\sigma\) is its nontrivial automorphism and \(\tilde\sigma\in G_{\Q}\)
a lift of \(\sigma\). For an elliptic curve \(E\subset A_K\) defined over \(K\),
\(E^\sigma=\tilde\sigma(E)\subset A_K\) is its conjugate; it is defined over \(K\) and does not depend
on the lift. We use the following results; the labels (I1)--(I4), (I7) and (I9) are those of
\cite[Section~1]{P2210}, and (I10) and (I11) are added here.

\begin{itemize}
\item \textbf{(I1) Reduction} \cite[(I1)]{P2210}, from Serre--Tate \cite{ST}. Let \(K'\) be a finite
extension of \(\Q_p\) with residue field \(k\) and \(A\) an abelian variety over \(K'\).
(a) Good reduction is preserved by \(K'\)-isogenies. (b) If \(A\) has good reduction
\(\widetilde A\) and \(m\) is prime to \(p\), then \(A(K')[m]\) injects into \(\widetilde A(k)\).
(c) If \(A\) has potentially good reduction, then it has good reduction over a finite extension of
\(K'\) with residue field \(k\). (d) An abelian surface \(B\) over \(\F_q\) has
\(\#B(\F_q)\le(1+\sqrt q)^4\).

\item \textbf{(I2) The case distinction} \cite[(I2)]{P2210}. A geometrically split abelian surface
\(A/\Q\) is of exactly one of the types \textbf{(N)} (\(A_{\bar\Q}\) is isogenous to \(E_1\times
E_2\) with \(E_1\), \(E_2\) not isogenous), \textbf{(CM)} (\(A_{\bar\Q}\) is isogenous to \(E^2\)
with \(E\) a CM curve) and \textbf{(S)} (\(A_{\bar\Q}\) is isogenous to \(E^2\) with \(E\) without
CM). In case (S), \(L\) is the smallest field over which all geometric endomorphisms are defined,
\(\Gamma=\Gal(L/\Q)\), and \(\Gamma\ne1\) if \(A\) is simple over \(\Q\).

\item \textbf{(I3)} \cite[Lemma~2]{P63}. Let \(X\) be an abelian surface over a number field
\(k\), \(C\subset X\) an elliptic curve defined over \(k\), and \(Q\in X(k)\) of prime order \(q\).
Then either \(Q\in C(k)\), or the image of \(Q\) in \((X/C)(k)\) has order \(q\).

\item \textbf{(I4)} \cite[Lemma~3]{P63}. In case (S), every elliptic curve in \(A_{\bar\Q}\) is
defined over \(L\).

\item \textbf{(I7) Weil restriction} \cite[(I7)]{P2210}. For an elliptic curve \(E\) over a
quadratic field \(K\), \((\Res_{K/\Q}E)(\Q)=E(K)\) and \((\Res_{K/\Q}E)_K\cong E\times E^\sigma\). If
\(E\) is the base change of an elliptic curve \(E_0/\Q\), then \(\Res_{K/\Q}E\) is isogenous over
\(\Q\) to \(E_0\times E_0^{(K)}\).

\item \textbf{(I9) Kubert's normal form} \cite[(I9)]{P2210}, from \cite[Table~3]{Ku}. For \(t\) in a
field \(k\) of characteristic \(0\) put
\[c=-\frac{t(t-1)(2t-1)}{t^2-3t+1},\qquad b=\frac{t^3(t-1)(2t-1)}{(t^2-3t+1)^2},\qquad
E_t:\ y^2+(1-c)xy-by=x^3-bx^2 .\]
If \(E_t\) is an elliptic curve, \((0,0)\) has order \(10\), and every pair \((E',P)\) of an
elliptic curve over \(k\) and a point \(P\in E'(k)\) of order \(10\) is isomorphic over \(k\) to
\((E_t,(0,0))\) for some such \(t\in k\).

\item \textbf{(I10) Torsion over quadratic fields} \cite[Theorem~(0.1)]{KM}, in the form quoted in
\cite[(I1)]{P2214}. Let \(N\) be a product of powers of \(2,3,5,7,11,13\) and \(m\mid N\). If
\(X_1(m,N)\) has positive genus and \((m,N)\) is not one of \((1,11)\), \((1,13)\), \((1,14)\), \((1,15)\),
\((1,16)\), \((1,18)\), \((2,10)\), \((2,12)\), then \(Y_1(m,N)(k)=\emptyset\) for every quadratic field \(k\). The curve \(X_1(20)\) has
genus \(3\) (certificate, part K). Hence \emph{no elliptic curve over a quadratic field has a point of
order \(20\).}

\item \textbf{(I11)} \cite[Section~3]{P63}, \cite[Section~5]{P2210}. If \(A\) is of type (N) and
simple over \(\Q\), the two elliptic curves of \(A_{\bar\Q}\) are defined over a quadratic field
\(K\) and exchanged by \(\sigma\). An elliptic curve \(E\subset A_K\) with \(E^\sigma=E\) is stable
under \(G_{\Q}\), hence an elliptic curve in \(A\) defined over \(\Q\) \cite[Section~2]{P2214}.
\end{itemize}

From \cite{P2210} we use the following results, with its numbering.
\begin{itemize}
\item \textbf{Lemma 1.} No abelian surface over \(\F_3\) has exactly \(40\) points.

\item \textbf{Proposition 2.} If \(A\) is of type (S), simple over \(\Q\), and \(A(\Q)[5]\ne0\),
then \(\Gamma\) is cyclic of order \(2\) or \(4\).

\item \textbf{Proposition A.} If \(A_{\bar\Q}\) is isogenous to the square of an elliptic curve
without CM and the field of definition of the geometric endomorphisms of \(A\) is a cyclic quartic
extension of \(\Q\), then \(\#A(\Q)[2^\infty]\le4\).
\end{itemize}
Neither Lemma~1, nor Proposition~2, nor Proposition~A uses a principal polarization
\cite[Remark~3]{P2210}. Proposition~2 rests on the residual-representation lemmas of
\cite[Section~6]{P63}, which use a polarization of \(A\) over \(\Q\) of arbitrary degree; every abelian
variety over \(\Q\) has one.

The following results of the quadratic case of \cite[Section~5]{P2210} concern an abelian surface
\(A/\Q\) that is simple over \(\Q\), a quadratic field \(K\), and an elliptic curve \(E\subset A_K\)
defined over \(K\) with \(E^\sigma\neq E\). In \cite{P2210} they are stated under the standing
assumption of that section. The proof of Lemma~5 uses only these hypotheses \cite[Remark~2]{P2210},
and Lemmas~6, 7 and Proposition~8 concern elliptic curves over \(K\) only.
\begin{itemize}
\item \textbf{Lemma 5 (Weil restriction).} (a) \(A\) is isogenous over \(\Q\) to
\(\Res_{K/\Q}E'\) for every elliptic curve \(E'\) over \(K\) that is \(K\)-isogenous to \(E\).
(b) If \(p\) is an odd prime, inert in \(K\), at which \(A\) has good reduction, then every such
\(E'\) has good reduction at \(pO_K\), and \(\#\widetilde A(\F_p)=\#\widetilde{E'}(O_K/pO_K)\).

\item \textbf{Lemma 6 (a Chebotarev criterion).} Let \(K=\Q(\sqrt d)\) and
\(E:y^2=(x-e_1)(x-e_2)(x-e_3)\) with \(e_i\in K\) distinct, \(N_{ij}=N_{K/\Q}(e_i-e_j)\). If
\(8\mid\#\widetilde E(O_K/pO_K)\) for all but finitely many primes \(p\) inert in \(K\), then for a
suitable numbering \(N_{ij}\in\Q^{\times2}\), and \(N_{ik}\in\Q^{\times2}\) or
\(N_{jk}\in\Q^{\times2}\) or \(dN_{ik}N_{jk}\in\Q^{\times2}\).

\item \textbf{Lemma 7 (\(\Z/2\times\Z/10\) over quadratic fields).} Let \(\mathcal X\) be the elliptic
curve \(y^2=x^3-2x^2+1\) over \(\Q\), with origin the point at infinity \(O\) (Cremona \(20a2\)), let
\(E_t\) be as in (I9), and put
\[
e_1=-\frac{t^3(t-1)}{t^2-3t+1},\qquad
B_t=\frac{-8t^4+12t^3-4t+1}{(t^2-3t+1)^2},
\]
\[
\rho=\frac{(2t-1)^2s}{(t^2-3t+1)^2},\qquad
e_{2,3}=\frac{-B_t\pm\rho}{8}.
\]
(a) \(\mathcal X\) has rank \(0\) and \(\mathcal X(\Q)=\{O,(0,\pm1),(1,0),(2,\pm1)\}\). (b) Let \(K\) be a
quadratic field and \(E'\) an elliptic curve over \(K\) with \(E'(K)\supseteq\Z/2\times\Z/10\). Then
\(E'\cong E_t\) over \(K\) for some \(t\in K\) such that \(s^2=(2t-1)(4t^2-2t-1)\) for some \(s\in K\),
and \(E'\) is isomorphic over \(K\) to \(y^2=(x-e_1)(x-e_2)(x-e_3)\). The point \(Q=(2t,s)\in\mathcal X(K)\)
does not lie in \(\mathcal X(\Q)\). (c) Put \(R=Q+Q^\sigma\in\mathcal X(\Q)\). If \(R=O\), then \(t\in\Q\).
If \(R\neq O\), then the line through \(Q\) and \(Q^\sigma\) passes through \(-R\) and has a slope
\(m\in\Q\); the \(x\)-coordinates of \(Q\) and \(Q^\sigma\) are the roots of the monic quadratic
\[q_{R,m}(X)=\frac{X^3-2X^2+1-\bigl(y(-R)+m(X-x(-R))\bigr)^2}{X-x(-R)}\in\Q[X],\]
which is irreducible over \(\Q\), and \(K=\Q\bigl(\sqrt{\disc q_{R,m}}\bigr)\).

\item \textbf{Proposition 8.} Let \(R\in\mathcal X(\Q)\setminus\{O\}\) and \(m\in\Q\) be such that
\(q_{R,m}\) is irreducible over \(\Q\). Let \(x_Q\) be a root of \(q_{R,m}\), \(K=\Q(x_Q)\), \(t=x_Q/2\) and
\(s=y(-R)+m(x_Q-x(-R))\), so that \(Q=(x_Q,s)\in\mathcal X(K)\), and suppose that \(E_t\) is an
elliptic curve. Then the conclusion of Lemma~6 fails for \(y^2=(x-e_1)(x-e_2)(x-e_3)\) over \(K\), with
\(e_1,e_2,e_3\) as above.
\end{itemize}
The computational inputs of Lemma~1, Lemma~7 and Proposition~8 are certified by
\texttt{qsimple\_2210\_checks.m} of \cite{P2210}.

\section{The CM case and endomorphism fields that are not quadratic}\label{other}

\textbf{Type (CM).} Let \(A\) be of type (CM) with a point \(P_5\in A(\Q)\) of order \(5\), and suppose
that \(\#A(\Q)[2^\infty]=2^k\ge8\). As in \cite[Section~2]{P2210}, \(A\) has potentially good reduction
at \(3\), and by (I1)(c) good reduction \(\widetilde A\) over a finite extension of \(\Q_3\) with
residue field \(\F_3\). The subgroup of \(A(\Q)\) generated by \(P_5\) and \(A(\Q)[2^\infty]\) has order
\(5\cdot2^k\), prime to \(3\), so by (I1)(b) \(5\cdot2^k\) divides \(\#\widetilde A(\F_3)\), which is at
most \((1+\sqrt3)^4<56\) by (I1)(d). Since \(5\cdot16=80>56\), this forces \(2^k=8\) and
\(\#\widetilde A(\F_3)=40\), contradicting Lemma~1. So \(\#A(\Q)[2^\infty]\le4\). This case uses neither
simplicity over \(\Q\) nor a polarization.

\textbf{Type (S).} Let \(A\) be of type (S), simple over \(\Q\), with a point of order \(5\). By
Proposition~2, \(\Gamma\cong C_2\) or \(\Gamma\cong C_4\). If \(\Gamma\cong C_4\), then
\(\#A(\Q)[2^\infty]\le4\) by Proposition~A.

It remains to treat type (S) with \(\Gamma\cong C_2\) and type (N).

\section{Moving \(2\)-primary torsion to an elliptic curve}\label{lemma}

\begin{statement}{Lemma 9.}\label{lemma-4prime}
Let \(A/\Q\) be an abelian surface, \(K\) a quadratic field and \(E\subset A_K\) an elliptic curve
defined over \(K\) with \(E^\sigma\neq E\). Put \(Y=A_K/E^\sigma\), and suppose that \(A(\Q)\) has a point
of order \(5\). Then:
\begin{enumerate}
\item[(a)] \(E(K)[2^\infty]\) and \(Y(K)[2^\infty]\) are killed by \(2\);
\item[(b)] \(A(\Q)[2^\infty]\) is killed by \(4\) and has order at most \(16\); in particular
\(A(\Q)\) has no point of order \(8\);
\item[(c)] if \(\#A(\Q)[2^\infty]\ge8\), there is an elliptic curve \(E'\) over \(K\), \(K\)-isogenous
to \(E\), such that \(E'(K)\supseteq\Z/2\times\Z/10\).
\end{enumerate}
\end{statement}

\emph{Proof.} \emph{Step 1: an exact sequence} (as in \cite[Proposition~2]{P2214}). The curves \(E\)
and \(E^\sigma\) are distinct elliptic curves in \(A_K\), so \(E\cap E^\sigma\) is finite. Let
\(\pi:A_K\to Y\) be the quotient map, defined over \(K\). Its restriction to \(E\) has the finite
kernel \(E\cap E^\sigma\), so it is a \(K\)-isogeny \(E\to Y\). If \(P\in A(\Q)\) and \(\pi(P)=0\),
then \(P\in E^\sigma\), hence \(P=\tilde\sigma(P)\in\tilde\sigma(E^\sigma)=\tilde\sigma^2(E)=E\), as
\(\tilde\sigma^2\in G_K\); so \(P\in E(\bar\Q)\cap A(K)=E(K)\). Therefore
\[
0\to A(\Q)[2^\infty]\cap E(K)\to A(\Q)[2^\infty]\xrightarrow{\ \pi\ }Y(K)[2^\infty]
\]
is exact, and \(\#A(\Q)[2^\infty]\le\#E(K)[2^\infty]\cdot\#Y(K)[2^\infty]\).

\emph{Step 2: a point of order \(5\).} Let \(P_5\in A(\Q)\) have order \(5\). By (I3) over \(K\),
applied to \(E^\sigma\subset A_K\), either \(P_5\in E^\sigma\), and then \(P_5\in E(K)\) as in Step 1,
or \(\pi(P_5)\in Y(K)\) has order \(5\). So for some \(C'\in\{E,Y\}\) there is a point
\(P\in C'(K)\) of order \(5\).

\emph{Step 3: combining} (as in \cite[Lemma~1]{P63}, with a group in place of a point). Let
\(C\in\{E,Y\}\) and \(T=C(K)[2^\infty]\). If \(C'=C\), put \(E'=C\), let \(a\) be the identity of
\(C\), and put \(T'=T\) and \(P'=P\). Otherwise let \(f:C\to C'\) be
the \(K\)-isogeny \(\pi|_E\) or its dual. Let \(H\) be the \(5\)-primary part of \(\ker f\); it is a
characteristic subgroup of the \(G_K\)-stable group \(\ker f\), hence \(G_K\)-stable. Put
\(E'=C/H\), let \(a:C\to E'\) be the quotient map and write \(f=b\circ a\) with \(b:E'\to C'\) defined
over \(K\). The degree of \(a\) is a power of \(5\), so \(a\) is injective on \(C[2^\infty]\) and
\(a(T)\subseteq E'(K)\) is isomorphic to \(T\). The degree of \(b\) is prime to \(5\); since \(b\circ\hat
b=[\deg b]\) is injective on \(C'[5]\), so is the dual isogeny \(\hat b:C'\to E'\), which is defined
over \(K\), and \(\hat b(P)\in E'(K)\) has order \(5\); put \(T'=a(T)\) and \(P'=\hat b(P)\). In both cases \(E'\) is \(K\)-isogenous to
\(E\), and
\[
E'(K)\supseteq T'\oplus\langle P'\rangle,\qquad T'\cong T,\quad P'\text{ of order }5 .
\]

\emph{(a)} If \(T\) had an element of order \(4\), then \(E'(K)\) would have a point of order \(20\),
contradicting (I10). So \(T\) is killed by \(2\), for \(C=E\) and for \(C=Y\).

\emph{(b)} By (a) and Step 1, \(\#A(\Q)[2^\infty]\le4\cdot4\). For \(x\in A(\Q)[2^\infty]\) we have
\(2\pi(x)=0\) by (a), so \(2x\in A(\Q)[2^\infty]\cap\ker\pi\subseteq E(K)[2^\infty]\), and \(4x=0\) by
(a).

\emph{(c)} If \(\#A(\Q)[2^\infty]\ge8\), Step 1 and (a) give \(C\in\{E,Y\}\) with \(\#T\ge4\) and \(T\)
killed by \(2\), so \(T=C[2]\). With \(E'\) and \(a\) as in Step 3, \(T'=a(T)\) is a subgroup of \(E'[2]\) of order
\(4\), so \(T'=E'[2]\). Hence \(E'[2]\subseteq E'(K)\) and \(E'(K)\supseteq\Z/2\times\Z/2\times\Z/5\cong
\Z/2\times\Z/10\). \qed

\medskip
Part (b) already excludes rational \(2\)-primary subgroups containing \(\Z/8\), without any Frobenius
argument; in particular it excludes \([40]\) in the cases of this section. The \(2\)-primary groups
\(\Z/2\times\Z/4\), \((\Z/2)^3\), \((\Z/4)^2\), \((\Z/2)^2\times\Z/4\) and \((\Z/2)^4\) need part (c) and
Section~\ref{quadratic}.

\section{The quadratic case}\label{quadratic}

\begin{statement}{Theorem Q.}\label{thm-Qprime}
Let \(A/\Q\) be an abelian surface that is simple over \(\Q\). Suppose that there are a quadratic
field \(K\) and an elliptic curve \(E\subset A_K\) defined over \(K\) with \(E^\sigma\neq E\). If
\(A(\Q)\) has a point of order \(5\), then \(\#A(\Q)[2^\infty]\le4\).
\end{statement}

\emph{Proof.} Suppose that \(\#A(\Q)[2^\infty]\ge8\). By
\hyperref[lemma-4prime]{Lemma~9(c)} there is an elliptic curve \(E'\) over \(K\), \(K\)-isogenous
to \(E\), with \(E'(K)\supseteq\Z/2\times\Z/10\). By Lemma~5(a), \(A\) is isogenous over \(\Q\) to
\(\Res_{K/\Q}E'\). Let \(p\) be an odd prime, inert in \(K\), at which \(A\) has good reduction. By
(I1)(b), \(A(\Q)[2^\infty]\) injects into \(\widetilde A(\F_p)\), so
\(8\mid\#\widetilde A(\F_p)=\#\widetilde{E'}(O_K/pO_K)\) by Lemma~5(b). This holds for all but
finitely many primes \(p\) inert in \(K\). Let \(E'\cong E_t\) and \(Q\in\mathcal X(K)\) be as in
Lemma~7.
\begin{itemize}
\item If \(Q+Q^\sigma=O\), then \(t\in\Q\), and \(E'\) is the base change of the elliptic curve
\(E_t\) over \(\Q\). By (I7), \(\Res_{K/\Q}E'\) is isogenous over \(\Q\) to \(E_t\times E_t^{(K)}\),
so \(A\) is not simple over \(\Q\). This is a contradiction.
\item Otherwise \(R=Q+Q^\sigma\neq O\). By Lemma~7(c), \(q_{R,m}\) is irreducible over \(\Q\), with
roots \(x(Q)\) and \(x(Q^\sigma)\), and \(K=\Q(x(Q))\) because \(x(Q)\notin\Q\). (If \(x(Q)\in\Q\), then
\(Q^\sigma\neq Q\) has the same \(x\)-coordinate, so \(Q^\sigma=-Q\) and \(R=O\).) As \(Q=(2t,s)\) lies on
the line through \(-R\) with slope \(m\), we have \(t=x(Q)/2\) and \(s=y(-R)+m(x(Q)-x(-R))\), and
\(E_t\cong E'\) is an elliptic curve. So the hypotheses of Proposition~8 hold with \(x_Q=x(Q)\). The
curve \(E'\) is isomorphic over \(K\) to \(y^2=(x-e_1)(x-e_2)(x-e_3)\), so the two curves have
isomorphic reductions at all but finitely many primes of \(K\), and Lemma~6 applies to
\(y^2=(x-e_1)(x-e_2)(x-e_3)\). Its conclusion contradicts Proposition~8.
\end{itemize}
Hence \(\#A(\Q)[2^\infty]\le4\). \qed

\medskip
\textbf{Type (S) with \(\Gamma\cong C_2\), and type (N).} Let \(A\) be simple over \(\Q\). If \(A\) is
of type (S) with \(\Gamma\cong C_2\), put \(K=L\); by (I4) every elliptic curve \(E\subset A_{\bar\Q}\)
is defined over \(K\). If \(A\) is of type (N), let \(K\) be the quadratic field of (I11) and \(E\) one
of the two elliptic curves of \(A_{\bar\Q}\). In both cases \(E^\sigma\ne E\) by (I11), since \(A\) is
simple over \(\Q\). So Theorem~Q applies.

\section{Proof of the Theorem and consequences for the census}\label{census}

\emph{Proof of the Theorem.} Let \(A/\Q\) be simple over \(\Q\) and geometrically split, with a
rational point of order \(5\). By (I2), \(A\) is of type (N), (CM) or (S). In type (CM),
\(\#A(\Q)[2^\infty]\le4\) by Section~\ref{other}. In type (S), \(\Gamma\cong C_2\) or \(C_4\) by
Proposition~2; for \(\Gamma\cong C_4\) the bound holds by Proposition~A, and for \(\Gamma\cong C_2\) by
Theorem~Q. In type (N) the bound holds by Theorem~Q. \qed

\medskip
\emph{Proof of the Corollary.} The Jacobian of a curve of genus \(2\) over \(\Q\) is an abelian surface.
The groups \([40]\) and \([2,20]\) have a point of order \(5\) and \(2\)-primary parts \(\Z/8\) and
\(\Z/2\times\Z/4\) of order \(8\); the groups \([2,2,10]\), \([2,2,20]\) and \([2,2,2,10]\) have
\(2\)-primary parts of orders \(8\), \(16\) and \(16\). \qed

\medskip
On the census, \([40]\) and \([2,20]\) become impossible in the class ``geometrically split, simple
over \(\Q\)''. Both have examples split over \(\Q\) (the census curves 40.b and 2.20.b), so their
combined geometrically split cells keep their grades. For \([2,2,10]\), \([2,2,20]\) and
\([2,2,2,10]\) the census entries already cite \cite{P2210}; the theorem removes its principal
polarization hypothesis. In the class ``geometrically split, simple over \(\Q\)'', the group \([20]\)
(\(2\)-primary part \(\Z/4\)) is not decided by this argument.

\section{Remarks}\label{remarks}

\begin{enumerate}
\item \emph{Where the hypotheses are used.} Simplicity over \(\Q\) gives \(E^\sigma\neq E\) in the
quadratic case, excludes the base-change alternative in Theorem~Q, and enters
Proposition~2. It is needed: the census lists examples split over \(\Q\) of \([40]\), \([2,20]\) and
\([2,2,10]\). The point of order \(5\) is also needed: the census curves 2.2.8.c and 2.2.2.4.c are
geometrically split and simple over \(\Q\), with rational torsion subgroups \([2,2,8]\) and
\([2,2,2,4]\), whose \(2\)-primary parts have order \(32\). In the quadratic case the rational point of
order \(5\) of \(A\) is used only in Step~2 of Lemma~9, to put a \(K\)-point of order \(5\) on \(E\) or
\(Y\).

\item \emph{The polarization in the earlier theorem.} By \cite[Remark~3]{P2210}, the principal
polarization was used there only in Lemma~4. Rational \(2\)-rank at least \(3\) gives
\(\#A(\Q)[2^\infty]\ge8\). So in the last paragraph of \cite[Section~5]{P2210}, Lemma~9(c) can
replace Lemma~4 and Lemma~5(c) together, with Lemma~5(a) applied to the original curve \(E\) (the
curve \(E'\) of Lemma~9 need not lie in \(A_K\)); the argument then becomes the proof of Theorem~Q.
Hence the theorem of \cite{P2210} holds for every abelian surface over \(\Q\)
that is simple over \(\Q\) and geometrically split; it is also a special case of the theorem above.

\item \emph{Isogenies and \(\Q\)-curves.} Rational \(2\)-primary torsion is not an isogeny invariant,
and \(\Q\)-isogenies of a Weil restriction \(\Res_{K/\Q}E'\) of a \(\Q\)-curve, for instance
\((2,2)\)-isogenies, can enlarge it. Lemma~9 bounds the \(2\)-primary torsion of every surface in
the \(\Q\)-isogeny class by that of two curves in the \(K\)-isogeny class of \(E\). If
\(E'(K)\supseteq\Z/2\times\Z/10\) and \(E'\) is not a base change, Proposition~8 and Lemma~6 give
infinitely many inert \(p\) with \(8\nmid\#\widetilde{E'}(\F_{p^2})=\#\widetilde A(\F_p)\), for every
\(A\) isogenous to \(\Res_{K/\Q}E'\). For example, the two elliptic curves over quadratic fields in the
LMFDB with torsion \(\Z/2\times\Z/10\) that are not base changes, \(2.2.13.1\text{-}36.1\text{-}a3\) and
\(2.2.17.1\text{-}36.1\text{-}a4\), have \(\#\widetilde E(\F_{121})=100\) (certificate, part S).

\item \emph{Sharpness.} Let \(E\) be the curve \(2.2.13.1\text{-}36.1\text{-}a3\),
\(y^2+(1+w)xy+(1+w)y=x^3+(-115+48w)x+(218-96w)\) over \(K=\Q(w)\), \(w^2-w-3=0\), and
\(A=\Res_{K/\Q}E\). The certificate (part S) checks that \(E(K)_{\mathrm{tors}}\cong\Z/2\times\Z/10\),
that \(j(E)=476379541/236196\) is not integral (so \(E\) has no CM), that \(E\cong E^\sigma\) over
\(K\), and that the Frobenius polynomial \(T^4-10T^2+49\) of \(A\) at the inert prime \(7\) is
irreducible. Hence \(A\) is simple over \(\Q\). As \(\End(E_{\bar\Q})=\Z\) and \(E\cong E^\sigma\) over
\(K\), every geometric endomorphism of \(A\) is defined over \(K\), so \(L\subseteq K\), and \(L\neq\Q\)
by (I2); thus \(A\) is of type (S) with \(\Gamma\cong C_2\). It is
principally polarized (the product polarization of \(E\times E^\sigma\) descends), and
\(A(\Q)_{\mathrm{tors}}=E(K)_{\mathrm{tors}}\cong\Z/2\times\Z/10\) by (I7). This \(A\) is not a Jacobian:
as \(\End(E_{\bar\Q})=\Z\), the principal polarizations of \(A_{\bar\Q}\cong E^2\) correspond to
positive definite unimodular binary quadratic forms over \(\Z\), each equivalent to \(x^2+y^2\), the
product polarization.

The bound is also attained by a Jacobian. Gluing \(E\) and \(E^\sigma\) along
\(\psi(\alpha_i,0)=(\sigma(\alpha_i),0)\), where \(\alpha_1,\alpha_2,\alpha_3\in K\) are the
\(x\)-coordinates of the points of order \(2\) of \(E\) \cite[Proposition~4]{HLP}, gives a curve over
\(K\) that descends to the genus-\(2\) curve
\[
C:\ y^2=63x^6+6x^5+184x^4+184x^2-6x+63
\]
over \(\Q\), the curve \(6084.d1\) of the LMFDB table \texttt{g2c\_curves\_new}. The certificate (part J)
checks that \(J(\Q)_{\mathrm{tors}}\cong\Z/2\times\Z/10\) for \(J=\mathrm{Jac}(C)\); that the
\(L\)-polynomial \(49T^4-10T^2+1\) of \(C\) at \(7\) is irreducible, so \(J\) is simple over \(\Q\); and
that over \(K\), \(C\) is isomorphic to the curve \(y^2=h(X)\) of \cite[Proposition~4]{HLP} with \(h\)
even, whose two elliptic quotients \((X,y)\mapsto(X^2,y)\) and \((X,y)\mapsto(1/X^2,y/X^3)\) are
\(K\)-isomorphic to \(E\) and to \(E^\sigma\cong E\); it also checks the non-degeneracy condition of
\cite[Proposition~4]{HLP}. The two quotient maps are defined over \(K\), so \(J_K\) is
\(K\)-isogenous to \(E^2\) and, as \(\End(E_{\bar\Q})=\Z\), every geometric endomorphism of \(J\) is
defined over \(K\). Since \(J\) is simple over \(\Q\), it is of type (S) with \(\Gamma\cong C_2\). The census lists \([2,10]\) in the class ``geometrically split, simple over \(\Q\)'' through
2.10.b (LMFDB 256.a.512.1), whose endomorphism field is cyclic quartic. In the LMFDB table
\texttt{g2c\_curves\_new}, five curves with torsion \([2,10]\) are simple over \(\Q\) and geometrically
split: \(6084.d1\), with Sato--Tate group \(E_2\), and \(256.a1\) (that is, 256.a.512.1), \(16900.i1\),
\(28900.u1\) and \(57600.gw1\), with Sato--Tate group \(E_4\); in the table \texttt{g2c\_curves} only
256.a.512.1 occurs.

\item \emph{Validation (not used in the proof).} Every curve in the class ``geometrically split,
simple over \(\Q\)'' with a rational point of order \(5\), in the census and in both LMFDB genus-\(2\)
tables, has \(2\)-primary torsion of order at most \(4\). Among the \(2795\) LMFDB elliptic curves \(E\)
over quadratic fields with \(5\mid\#E(K)_{\mathrm{tors}}\), each of the \(2108\) for which
\(\Res_{K/\Q}E\) is certified simple over \(\Q\) has a good odd prime \(p\le23\) with
\(8\nmid\#\widetilde{\Res_{K/\Q}E}(\F_p)\). Across the
\(2403\) isogeny classes of LMFDB elliptic curves over quadratic fields that contain a curve with a
point of order \(5\), no curve has a point of order \(4\), as (I10) predicts.

\item \emph{Two derivations.} The argument of Sections~\ref{lemma} and~\ref{quadratic} was found
independently in two separate analyses, of type (N) and of type (S) with \(\Gamma\cong C_2\); each
checked the other.
\end{enumerate}

\section{Computations}\label{computations}

All scripts run in Magma V2.30-1 and consist of assertions.
\begin{itemize}
\item \texttt{qsimple\_40\_checks.m}, which runs in under a second. Its log
\texttt{qsimple\_40\_checks.log} ends with \texttt{ALL\_T5SQ\_CHECKS\_PASSED}. Part K:
\(X_1(20)\) has genus \(3\) and \((1,20)\) is not among the exceptions of (I10). Parts S and J: the facts of
Remarks~3 and~4. The SHA-256 hash of the script is
\texttt{2a06ab43a45d3617\allowbreak f00cc3ffc82b5e88\allowbreak e5ebaf7c255e5805\allowbreak 1c8bd966ace1fce2}, and that of the log is
\texttt{a7ff82365c9fcddc\allowbreak f8b598df50a23a56\allowbreak abf8afa71192bced\allowbreak c1778c1a9dc43535}.
\item The published certificate \texttt{qsimple\_2210\_checks.m} of \cite{P2210} certifies Lemma~1,
Lemma~7 and Proposition~8. Rerun unchanged, it gives a log byte-identical to the published one.
\end{itemize}

\begin{thebibliography}{P2214}

\bibitem[HLP]{HLP}
E. W. Howe, F. Leprévost and B. Poonen, \emph{Large torsion subgroups of split Jacobians of curves of
genus two or three}, Forum Math. \textbf{12} (2000), 315--364, Proposition~4.

\bibitem[KM]{KM}
M. A. Kenku and F. Momose, \emph{Torsion points on elliptic curves defined over quadratic fields},
Nagoya Math. J. \textbf{109} (1988), 125--149, Theorem~(0.1).
\href{https://doi.org/10.1017/S0027763000002816}{doi:10.1017/S0027763000002816}.

\bibitem[Ku]{Ku}
D. S. Kubert, \emph{Universal bounds on the torsion of elliptic curves}, Proc. London Math. Soc. (3)
\textbf{33} (1976), 193--237, Table~3.
\href{https://doi.org/10.1112/plms/s3-33.2.193}{doi:10.1112/plms/s3-33.2.193}.

\bibitem[P63]{P63}
Census proof of 21 October 2026, \emph{An obstruction to rational \(63\)-torsion on geometrically
split abelian surfaces},
\url{https://genus-2-torsion.github.io/proofs/qsimple-63-obstruction.html}: Lemmas~1--3 and
Sections~3 and~6.

\bibitem[P2210]{P2210}
Census proof of 11 November 2026, \emph{Rational \(5\)-torsion and the \(2\)-rank of geometrically
split Jacobians},
\url{https://genus-2-torsion.github.io/proofs/qsimple-2210-obstruction.html}: Sections~1, 2 and 5,
Lemma~1, Propositions~2, A and 8, Lemmas~5--7, Remarks~2 and~3.

\bibitem[P2214]{P2214}
Census proof of 10 November 2026, \emph{Rational \(7\)-torsion and the \(2\)-rank of
geometrically split abelian surfaces},
\url{https://genus-2-torsion.github.io/proofs/qsimple-2214-obstruction.html}: (I1), Section~2 and
Proposition~2.

\bibitem[ST]{ST}
J.-P. Serre and J. Tate, \emph{Good reduction of abelian varieties}, Ann. of Math. (2)
\textbf{88} (1968), 492--517. \href{https://doi.org/10.2307/1970722}{doi:10.2307/1970722}.

\end{thebibliography}

\end{document}
